\ifdefined\gkver\else\def\gkver{student}\fi
\documentclass[\gkver]{gaokaozhenti}
\begin{document}

\chapter{2001年粤豫}

\begin{enumerate}
\item
%% number: 1
%% paperName: 2001年粤豫高考第 1 题 4 分
%% typeId: 2
%% type: 多选题
%% chapter: 机械振动
%% point: 单摆
%% method: 无
%% score: 4
%% degree: 850
%% duplicateId: 0
%% body:
一单摆做简谐振动，对摆球所经过的任何一点来说，相继两次通过该点时，摆球的\xzanswer{BD}
\fourchoices[answer=BD]
{速度必相同}
{加速度必相同}
{动量必相同}
{动能必相同}
%% answer: BD
%% memo:
\memoanswer{
本题原卷及材料中未提供详解，待补充。
}

\item
%% number: 6
%% paperName: 2001年粤豫高考第 6 题 4 分
%% typeId: 2
%% type: 多选题
%% chapter: 机械波
%% point: 波的多解性
%% method: 无
%% score: 4
%% degree: 450
%% duplicateId: 0
%% body:
一列简谐横波在图中 $x$ 轴上传播，a、b 是其中相距为 $0.3\Um$ 的两点。在某时刻，a 点质元正位于平衡位置向上运动，b 点质元恰好运动到下方最大位移处。已知横波的传播速度为 $60\Ums$，波长大于 $0.3\Um$。\xzanswer{AD}
\onepicture[w=7.5cm]{\includesvg{figs/06}}
\fourchoices[answer=AD]
{若该波沿 $x$ 轴负方向传播，则频率为 $150\UHz$}
{若该波沿 $x$ 轴负方向传播，则频率为 $100\UHz$}
{若该波沿 $x$ 轴正方向传播，则频率为 $75\UHz$}
{若该波沿 $x$ 轴正方向传播，则频率为 $50\UHz$}
%% answer: AD
%% memo:
\memoanswer{
本题原卷及材料中未提供详解，待补充。
}

\item
%% number: 10
%% paperName: 2001年粤豫高考第 10 题 4 分
%% typeId: 1
%% type: 单选题
%% chapter: 电场
%% point: 电势差与电场强度的关系
%% method: 无
%% score: 4
%% degree: 550
%% duplicateId: 0
%% body:
如图所示，平行板电容器经开关 K 与电池连接，a 处有一带电量非常小的点电荷。K 是闭合的，$\varphi_{a}$ 表示 a 点的电势，$F$ 表示点电荷受到的电场力。现将电容器的 B 板向下稍微移动，使两板间的距离增大，则\xzanswer{B}
\onepicture[w=4.8cm]{figs/10.png}
\fourchoices[answer=B]
{$\varphi_{a}$ 变大，$F$ 变大}
{$\varphi_{a}$ 变大，$F$ 变小}
{$\varphi_{a}$ 不变，$F$ 不变}
{$\varphi_{a}$ 不变，$F$ 变小}
%% answer: B
%% memo:
\memoanswer{
由题意知，B 板下移，板间距离 $d$ 变大，电容器电容变小；K 闭合，电容器电压不变，由 $E = \frac{U}{d}$ 知场强 $E$ 变小，由 $F = qE$ 知电场力 $F$ 变小；由 $U_{Aa} = Ed_{Aa} = \varphi_{A} - \varphi_{a}$（a 点电势等于 A 板的电势（不变）减去 a 点与 A 板的电势差），得到 $\varphi_{a}$ 大，故 B 正确。
}

\end{enumerate}

\end{document}
